% =====================================================================
%  defs.tex — Definicions, entorns i comandes reutilitzables
%  Matemàtiques Aplicades I · 1r Batxillerat PRO · Unitat Didàctica 2
%  Es carrega una sola vegada des de main.tex (després del preàmbul).
% =====================================================================

% ------- Colors de la col·lecció (sobris, pocs tons) -----------------
\definecolor{udnavy}{HTML}{1F3A5F}   % blau fosc (títols)
\definecolor{udblue}{HTML}{2E75B6}   % blau accent
\definecolor{udbg}{HTML}{F4F6F9}     % fons gris molt clar
\definecolor{udline}{HTML}{DFE3EA}   % línies suaus
\definecolor{udgreen}{HTML}{2F7D52}  % per a solucions
\definecolor{udgrey}{HTML}{5F5E5A}   % text secundari

% ------- Comptadors --------------------------------------------------
\newcounter{exemple}[section]
\renewcommand{\theexemple}{\thesection.\arabic{exemple}}
\newcounter{exresolt}[section]
\renewcommand{\theexresolt}{\thesection.\arabic{exresolt}}
\newcounter{exproposat}[section]
\renewcommand{\theexproposat}{\thesection.\arabic{exproposat}}

% =====================================================================
%  Entorn: EXEMPLE (concepte il·lustrat)
% =====================================================================
\newtcolorbox{exemple}[1][]{
  enhanced, breakable,
  colback=udbg, colframe=udblue,
  boxrule=0pt, leftrule=3pt,
  arc=2pt, left=10pt, right=10pt, top=8pt, bottom=8pt,
  coltitle=udnavy, fonttitle=\bfseries,
  title={\stepcounter{exemple}Exemple \theexemple\ifx\relax#1\relax\else\ — #1\fi},
}

% =====================================================================
%  Entorn: EXERCICI RESOLT
% =====================================================================
\newtcolorbox{exresolt}[1][]{
  enhanced, breakable,
  colback=white, colframe=udnavy,
  boxrule=0.6pt, arc=2pt,
  left=10pt, right=10pt, top=8pt, bottom=8pt,
  coltitle=white, fonttitle=\bfseries, colbacktitle=udnavy,
  attach boxed title to top left={yshift=-2mm, xshift=4mm},
  boxed title style={colframe=udnavy, arc=2pt},
  title={\stepcounter{exresolt}Exercici resolt \theexresolt\ifx\relax#1\relax\else\ — #1\fi},
}

% Etiqueta interna "Solució:" per als exercicis resolts
\newcommand{\solucio}{\par\medskip\noindent\textbf{\color{udgreen}Solució.}\ }
\newcommand{\resposta}[1]{\par\medskip\noindent
  \colorbox{udbg}{\parbox{\dimexpr\linewidth-2\fboxsep}{%
  \textbf{\color{udnavy}Resposta:} #1}}}

% =====================================================================
%  Entorn: EXERCICIS PER FER (llista numerada amb capçalera)
% =====================================================================
\newtcolorbox{exercicis}[1][Exercicis per fer]{
  enhanced, breakable,
  colback=white, colframe=udblue,
  boxrule=0pt, leftrule=3pt, toprule=0pt, bottomrule=0pt, rightrule=0pt,
  sharp corners,
  left=10pt, right=10pt, top=6pt, bottom=6pt,
  coltitle=udnavy, fonttitle=\bfseries,
  title={#1},
}

% Llista numerada per als exercicis (numeració global per secció)
\newlist{exlist}{enumerate}{1}
\setlist[exlist]{label=\textbf{\thesection.\arabic*}, ref=\thesection.\arabic*,
  leftmargin=2.6em, itemsep=4pt, topsep=4pt}

% =====================================================================
%  Comoditats tipogràfiques
% =====================================================================
% Vector d'idees clau / nota al marge conceptual
\newtcolorbox{idea}{
  enhanced, breakable,
  colback=udbg, colframe=udline,
  boxrule=0.6pt, arc=2pt,
  left=9pt, right=9pt, top=6pt, bottom=6pt,
}

% Encapçalament d'sessió (s'usa al començament de cada fitxer d'sessió)
\newcommand{\horatitol}[2]{%
  \section{#1}%
  \noindent{\color{udgrey}\itshape #2}\par\medskip
}

% Euros i multiplicació "social" coherents
\newcommand{\eur}{\,\text{\euro}}
\newcommand{\per}{\cdot}

% =====================================================================
%  Entorn: SOLUCIONS (full per al professorat)
%  Només resultats/respostes, sense resolució detallada.
% =====================================================================
\newtcolorbox{solucions}[1][]{
  enhanced, breakable,
  colback=udbg, colframe=udgreen,
  boxrule=0pt, leftrule=3pt,
  arc=2pt, left=10pt, right=10pt, top=8pt, bottom=8pt,
}

% Capçalera del full de solucions (banda visible "només professorat")
\newcommand{\fullsolucions}[1]{%
  \newpage
  \begin{center}
    \colorbox{udgreen}{\color{white}\bfseries\;Full del professorat — Solucions\;}
  \end{center}
  \subsection*{Solucions — #1}
  \addcontentsline{toc}{subsection}{Solucions — #1}
}

% Llista de solucions: una entrada per exercici proposat
\newlist{sollist}{enumerate}{1}
\setlist[sollist]{label=\textbf{\thesection.\arabic*}, leftmargin=2.6em,
  itemsep=3pt, topsep=3pt}

% =====================================================================
%  Gràfiques de funcions (pgfplots) — estil comú de la col·lecció
%  Requereix tikz + pgfplots carregats al preàmbul (main.tex / wrapper).
%  Es defineix aquí perquè usa els colors udgrey i udline ja existents.
% =====================================================================
\pgfplotsset{
  udaxis/.style={
    axis lines=middle,
    axis line style={-{Stealth[length=2mm]}, udgrey},
    tick style={udgrey},
    label style={font=\small\color{udgrey}},
    tick label style={font=\scriptsize\color{udgrey}},
    grid=both,
    major grid style={line width=0.2pt, draw=udline},
    minor tick num=0,
    every axis plot/.append style={line width=1.1pt},
    width=8.4cm, height=6cm,
  },
}

% Peu de figura breu i sobri (s'escriu dins d'un center, després del dibuix).
\newcommand{\figcap}[1]{\par\smallskip{\small\color{udgrey}\itshape #1}}
